\subsection{Frame Semantics}
\label{sec:frame-semantics}
\slidesrc{Other semantics: 2--18; Semantics (long version): 10}

\subsubsection{Word meaning revisited}
\slidesrc{Other semantics: 2--6}

What is the meaning of \emph{sold}?
\begin{itemize}
  \item \textbf{Distributional semantics}: all the contexts in which \emph{sold} occurs
        (\emph{\ldots was sold by\ldots}, \emph{\ldots sold me that piece of\ldots}). It can
        find similar words and contexts and generalise (e.g.\ with dimensionality
        reduction), but the representation is \textbf{monolithic}: a word vector has no
        internal structure, nothing in it says ``there is a seller and a buyer''.
  \item \textbf{Frame semantics}: the meaning is given by a \textbf{frame}, a
        stereotypical situation with \textbf{roles} (slots) to be filled.
\end{itemize}

\begin{nbdef}{Frame}
A \textbf{frame} is a structured representation of a stereotypical situation, made of
named \textbf{roles} (frame elements). A word \emph{evokes} the frame, and the sentence
fills its roles.
\end{nbdef}

\begin{nbexample}{a commercial transaction}
\emph{Cynthia sold the bike for \$200.}
\begin{center}
\begin{tabular}{@{}l l@{}}
\multicolumn{2}{@{}l}{\textbf{Commercial transaction}}\\
\midrule
\textsc{seller} & Cynthia \\
\textsc{buyer}  & ? \\
\textsc{goods}  & the bike \\
\textsc{price}  & \$200 \\
\end{tabular}
\end{center}
The \textsc{buyer} role stays empty: the sentence does not say who bought the bike, but
the frame tells us that \emph{someone} did.
\end{nbexample}

\subsubsection{Two properties of frames}
\slidesrc{Other semantics: 7--11}

Following Fillmore (1977) and Langacker (1987):

\begin{itemize}
  \item \textbf{Prototypical}: a frame does not need to handle all the cases. The frame
        of \emph{widow} is ``a woman marries one man, and the man dies''. What if a woman
        had three husbands, two of whom died? The frame does not say, and it does not
        need to: it describes the typical situation.
  \item \textbf{Profiling}: words evoking the same frame can \emph{highlight} different
        aspects of it.
        \begin{itemize}[nosep]
          \item \emph{sell} is seller-centric, \emph{buy} is buyer-centric:
                \emph{Cynthia sold the bike (to Bob)} vs.\ \emph{Bob bought the bike
                (from Cynthia)};
          \item \emph{rob} highlights the person, \emph{steal} the goods:
                \emph{Cynthia robbed Bob (of the bike)} vs.\ \emph{Cynthia stole the bike
                (from Bob)}.
        \end{itemize}
        The profiled role becomes the subject or object; the others are optional (in
        brackets).
\end{itemize}

\subsubsection{From syntax to semantics}
\slidesrc{Other semantics: 12--15}

The same situation, with \textsc{seller} Cynthia, \textsc{buyer} Bob, \textsc{goods} the
bike, \textsc{price} \$200, can be expressed by many \textbf{syntactic alternations},
with different verbs and different arguments in different positions:
\begin{itemize}[nosep]
  \item Cynthia sold the bike to Bob for \$200.
  \item The bike sold for \$200.
  \item Bob bought the bike from Cynthia.
  \item The bike was bought by Bob.
  \item The bike was bought for \$200.
  \item The bike was bought for \$200 by Bob.
\end{itemize}
The subject is the seller in the first sentence, the goods in the second, the buyer in
the third. \textbf{Goal}: map \emph{syntactic positions} to \emph{semantic roles}.

\subsubsection{Semantic role labelling}
\slidesrc{Other semantics: 16--18}

\begin{nbdef}{Semantic role labelling (SRL)}
Given a sentence, label each relevant phrase with its semantic role in the frame evoked
by the predicate:
\begin{center}
\begin{tabular}{@{}l ccccc@{}}
\textbf{Input:}  & Cynthia & sold & the bike & to Bob & for \$200 \\
\textbf{Output:} & \textsc{seller} & \textsc{predicate} & \textsc{goods} & \textsc{buyer} & \textsc{price}
\end{tabular}
\end{center}
Subtasks:
\begin{enumerate}[nosep]
  \item \textbf{frame identification}: find the \textsc{predicate} and the frame it
        evokes;
  \item \textbf{argument identification}: find the phrases filling \textsc{seller},
        \textsc{goods}, etc.
\end{enumerate}
\end{nbdef}

\begin{nbintuition}[SRL needs syntax]
\emph{Bob} is the \textsc{buyer} in ``Cynthia sold the bike \emph{to Bob}'' and in
``\emph{Bob} bought the bike'', but in different syntactic positions. SRL systems
therefore use the output of POS tagging and parsing
(Section~\ref{sec:sentence-level-syntax}): syntax is the input of semantics.
\end{nbintuition}

\textbf{Tools}: \textbf{FrameNet} (\url{https://framenet.icsi.berkeley.edu/}) is a
lexical database of English frames, with their roles and annotated example sentences;
see for example the frame \emph{Commerce\_buy}.
